% ANEXOS I
\setcounter{section}{0}

\newcommand{\VIUAnexoSection}[2]{%
  \refstepcounter{section}%
  \VIUSectionStar{#1}%
  \label{#2}%
}

\VIUAnexoSection{Anexo I}{ane:anexo_checklist_calidad}
\subsection*{ Checklist de evaluación de calidad}

\begin{longtblr}
[
 label={tab:anexo_checklist_calidad},
 caption={Checklist de evaluación de calidad adaptado de Kitchenham (2007), con aplicación condicionada por subtipo cuantitativo. Elaboración propia.}] {
  colspec = {X[0.8,j,m] X[3.5,j,m] X[1.2,j,m] X[2,j,m]},
  width = \textwidth
}
\textbf{Código} & \textbf{Pregunta operativa} & \textbf{Etapa} & \textbf{Aplicabilidad por subtipo} \\

\texttt{QA\_I01} & ¿Los objetivos del estudio están claramente definidos? & Diseño & General (Exp, Corr, Q-NS) \\
\texttt{QA\_I02} & ¿Las medidas definidas permiten responder la pregunta del estudio? & Diseño & Solo Exp \\
\texttt{QA\_I03} & ¿Se justifica el tamaño de muestra o número de observaciones? & Diseño & Solo Exp \\
\texttt{QA\_I04} & ¿La tecnología/sistema evaluado está claramente descrito? & Diseño & General (Exp, Corr, Q-NS) \\
\texttt{QA\_I05} & ¿Las variables clave se miden adecuadamente? & Ejecución & General (Exp, Corr, Q-NS) \\
\texttt{QA\_I06} & ¿Las medidas e indicadores están definidos de forma completa? & Ejecución & General (Exp, Corr, Q-NS) \\
\texttt{QA\_I07} & ¿Las medidas seleccionadas son relevantes para la pregunta del estudio? & Diseño & General (Exp, Corr, Q-NS) \\
\texttt{QA\_I08} & ¿El alcance del estudio es suficiente para sostener sus conclusiones? & Diseño & Exp y Q-NS \\
\texttt{QA\_I09} & ¿Los métodos de recolección de datos están descritos? & Ejecución & General (Exp, Corr, Q-NS) \\
\texttt{QA\_I10} & ¿Se explican los tipos y fuentes de datos utilizados? & Ejecución & General (Exp, Corr, Q-NS) \\
\texttt{QA\_I11} & ¿Las unidades de medida están claramente especificadas? & Ejecución & General (Exp, Corr, Q-NS) \\
\texttt{QA\_I12} & ¿Se reportan los datos básicos necesarios para interpretar resultados? & Ejecución & General (Exp, Corr, Q-NS) \\
\texttt{QA\_I13} & ¿Se describen los métodos estadísticos o analíticos empleados? & Análisis & General (Exp, Corr, Q-NS) \\
\texttt{QA\_I14} & ¿Se referencia el software/herramienta usada para el análisis? & Análisis & General (Exp, Corr, Q-NS) \\
\texttt{QA\_I15} & ¿Se justifica la elección de los métodos de análisis? & Análisis & General (Exp, Corr, Q-NS) \\
\texttt{QA\_I16} & ¿El propósito del análisis está claramente establecido? & Análisis & General (Exp, Corr, Q-NS) \\
\texttt{QA\_I17} & ¿Se describe el sistema de puntuación o criterios de evaluación del experimento? & Análisis & Solo Exp \\
\texttt{QA\_I18} & ¿Se controlan o discuten variables de confusión relevantes? & Análisis & General (Exp, Corr, Q-NS) \\
\texttt{QA\_I19} & ¿Existe consistencia numérica entre tablas, texto y resultados? & Análisis & General (Exp, Corr, Q-NS) \\
\texttt{QA\_I20} & ¿Se aborda el riesgo de sesgo de selección? & Análisis & Exp y Q-NS \\
\texttt{QA\_I21} & ¿El estudio responde explícitamente sus preguntas/objetivos? & Resultados y discusión & General (Exp, Corr, Q-NS) \\
\texttt{QA\_I22} & ¿Los hallazgos principales se interpretan con claridad? & Resultados y discusión & General (Exp, Corr, Q-NS) \\
\texttt{QA\_I23} & ¿Se reportan también hallazgos negativos o no favorables? & Resultados y discusión & General (Exp, Corr, Q-NS) \\
\texttt{QA\_I24} & ¿Se discute la significancia práctica de los resultados? & Resultados y discusión & General (Exp, Corr, Q-NS) \\
\texttt{QA\_I25} & ¿Se interpretan resultados nulos o no concluyentes? & Resultados y discusión & General (Exp, Corr, Q-NS) \\
\texttt{QA\_I26} & ¿Se discute si pudieron omitirse efectos relevantes? & Resultados y discusión & General (Exp, Corr, Q-NS) \\
\texttt{QA\_I27} & ¿Los resultados se comparan con estudios previos? & Resultados y discusión & General (Exp, Corr, Q-NS) \\
\texttt{QA\_I28} & ¿Se explica el aporte del estudio al estado del arte? & Resultados y discusión & General (Exp, Corr, Q-NS) \\
\texttt{QA\_I29} & ¿Se presentan implicaciones para práctica/sistemas reales? & Resultados y discusión & General (Exp, Corr, Q-NS) \\
\texttt{QA\_I30} & ¿Se explican consecuencias para validez/confiabilidad y trabajo futuro? & Resultados y discusión & General (Exp, Corr, Q-NS) \\

\end{longtblr}


\VIUAnexoSection{Anexo II}{ane:anexo_diccionario_extraccion}
\subsection*{ Diccionario operativo de la matriz de extracción}
%Los anexos también contienen información adicional que se considera relevante para justificar las conclusiones del trabajo, pero, por lo general, el autor de contenido del anexo es distinto al autor del trabajo. Suele ser un documento independiente del trabajo. Pueden ser tablas de datos, imágenes, etc. Es necesario incluir las referencias de los documentos de donde procedan.

\begin{longtblr}
[
 label={tab:diccionario_extraccion},
 caption={Diccionario operativo de campos de la matriz final de extracción de datos.}]{
  colspec = {X[0.8,j,m] X[2.5,j,m] X[4.0,j,m]},
  rowhead = 1,
  row{1} = {font=\bfseries},
  hlines
}
Bloque & Campo & Descripción operacional \\

Meta & \texttt{Article\_ID} & Identificador único del estudio en la revisión. \\
Meta & \texttt{Title} & Título del artículo. \\
Meta & \texttt{Year} & Año de publicación. \\
Meta & \texttt{Venue} & Revista o congreso de publicación. \\
Meta & \texttt{DOI\_or\_URL} & DOI o enlace trazable de la fuente. \\
Meta & \texttt{Study\_type\_quant} & Subtipo cuantitativo aplicado en evaluación metodológica. \\
Meta & \texttt{Country\_or\_region} & País o región del caso de estudio cuando es reportado. \\

D1 & \texttt{D1\_Pollutants\_inputs} & Contaminantes atmosféricos usados como entradas. \\
D1 & \texttt{D1\_Env\_variables\_inputs} & Variables ambientales no contaminantes usadas como entradas. \\
D1 & \texttt{D1\_Target\_output} & Variable objetivo reportada por el estudio. \\
D1 & \texttt{D1\_Reported\_input\_values} & Rangos o valores numéricos de variables de entrada reportados por el estudio. \\
D1 & \texttt{D1\_Sampling\_frequency\_values} & Frecuencia de muestreo o periodicidad de adquisición. \\
D1 & \texttt{D1\_Time\_window\_values} & Ventana temporal, periodo de captura o duración del dataset. \\
D1 & \texttt{D1\_Data\_source\_type} & Tipo de fuente de datos (sensores, datasets públicos, híbrido, etc.). \\
D1 & \texttt{D1\_Preprocessing} & Pasos de preprocesamiento aplicados antes del modelado. \\

D2 & \texttt{D2\_Fuzzy\_approach\_type} & Tipo de enfoque difuso implementado (Mamdani, Sugeno, ANFIS, etc.). \\
D2 & \texttt{D2\_Fuzzy\_role} & Rol del componente difuso en la solución (central, comparativo, complementario). \\
D2 & \texttt{D2\_Model\_size\_values} & Tamaño/estructura del modelo (entradas, reglas, capas, estados). \\
D2 & \texttt{D2\_Optimization\_values} & Parámetros o métodos de optimización/entrenamiento reportados. \\
D2 & \texttt{D2\_Inference\_purpose} & Propósito funcional de la inferencia difusa (estimación, clasificación, alerta, ruta). \\
D2 & \texttt{D2\_Rule\_base\_definition} & Forma de definición de la base de reglas. \\
D2 & \texttt{D2\_Membership\_functions} & Tipo de funciones de pertenencia y representación lingüística. \\
D2 & \texttt{D2\_Defuzzification\_method} & Método de defuzzificación reportado. \\
D2 & \texttt{D2\_Risk\_modeling\_approach} & Enfoque para traducir inferencia a riesgo o severidad. \\
D2 & \texttt{D2\_Interpretability\_elements} & Elementos explícitos de interpretabilidad (reglas, tablas, gráficos, etc.). \\

D3 & \texttt{D3\_Platform\_architecture} & Arquitectura de plataforma/sistema de monitoreo. \\
D3 & \texttt{D3\_Processing\_mode} & Modo de procesamiento (tiempo real, casi real, híbrido, offline). \\
D3 & \texttt{D3\_Communication\_and\_tech} & Tecnologías de comunicación, software y servicios usados. \\
D3 & \texttt{D3\_Visualization\_alerting} & Mecanismos de visualización y alertamiento implementados. \\
D3 & \texttt{D3\_Deployment\_context} & Contexto de despliegue (indoor/outdoor, urbano, industrial, etc.). \\
D3 & \texttt{D3\_Hardware\_and\_sensors} & Hardware y sensores utilizados en implementación/medición. \\
D3 & \texttt{D3\_Real\_time\_characteristics} & Características operativas de comportamiento en tiempo real. \\

Eval & \texttt{Eval\_Metrics\_reported} & Métricas de desempeño reportadas por el estudio. \\
Eval & \texttt{Eval\_Reported\_values} & Valores numéricos de las métricas reportadas. \\
Eval & \texttt{Eval\_Comparator\_values} & Resultados de comparación con baseline/modelos alternativos. \\
Eval & \texttt{Eval\_Improvement\_reported} & Mejoras cuantificadas o declaradas frente a referencias. \\
Eval & \texttt{Eval\_Dataset\_size\_values} & Tamaño del dataset, muestras o registros usados en evaluación. \\
Eval & \texttt{Eval\_Data\_split\_values} & Esquema de partición de datos para entrenamiento/validación/prueba. \\
Eval & \texttt{Eval\_Key\_results\_summary} & Síntesis técnica breve de resultados principales del estudio. \\

Synth & \texttt{Main\_contribution\_reported} & Contribución principal declarada para síntesis narrativa. \\
Synth & \texttt{Main\_limitation\_reported} & Limitación principal reportada para discusión crítica. \\
Synth & \texttt{Evidence\_D1\_pages} & Páginas de soporte para campos D1. \\
Synth & \texttt{Evidence\_D2\_pages} & Páginas de soporte para campos D2. \\
Synth & \texttt{Evidence\_D3\_pages} & Páginas de soporte para campos D3. \\
Synth & \texttt{Evidence\_Eval\_pages} & Páginas de soporte para resultados de evaluación. \\
Synth & \texttt{QA\_Quality\_band} & Banda de calidad metodológica del estudio (Alta/Media/Baja). \\
Synth & \texttt{QA\_Score\_percent} & Puntaje porcentual de calidad metodológica. \\
Synth & \texttt{Sensitivity\_set} & Etiqueta de inclusión para análisis de sensibilidad de síntesis. \\
\end{longtblr}



\clearpage
\VIUAnexoSection{Anexo III}{ane:d3_tabla_principal}
\subsection*{Tabla principal de evidencia cuantitativa}
\input{assets/tablas/03_finales/03_anexos_i/anexo_d3_tabla_principal_depurada.tex}

\clearpage
\VIUAnexoSection{Anexo IV}{ane:anexo_reglas_principales}
\subsection*{Base principal de 54 reglas del motor difuso}
\begin{longtblr}[
  label={tab:anexo_reglas_principales_54},
  caption={Base principal de 54 reglas del motor difuso implementado.}
]{
  colspec={X[1.0,l,m] X[1.1,l,m] X[1.1,l,m] X[1.7,l,m]},
  width=\linewidth
}
\toprule
AQI base & Concurrencia & Persistencia & Salida \\
\midrule
Good & Baja & Baja & Good \\
Good & Baja & Media & Good \\
Good & Baja & Alta & Moderate \\
Good & Media & Baja & Good \\
Good & Media & Media & Moderate \\
Good & Media & Alta & Moderate \\
Good & Alta & Baja & Moderate \\
Good & Alta & Media & Moderate \\
Good & Alta & Alta & Unhealthy for Sensitive Groups \\
Moderate & Baja & Baja & Moderate \\
Moderate & Baja & Media & Moderate \\
Moderate & Baja & Alta & Unhealthy for Sensitive Groups \\
Moderate & Media & Baja & Moderate \\
Moderate & Media & Media & Unhealthy for Sensitive Groups \\
Moderate & Media & Alta & Unhealthy for Sensitive Groups \\
Moderate & Alta & Baja & Unhealthy for Sensitive Groups \\
Moderate & Alta & Media & Unhealthy for Sensitive Groups \\
Moderate & Alta & Alta & Unhealthy \\
Unhealthy for Sensitive Groups & Baja & Baja & Unhealthy for Sensitive Groups \\
Unhealthy for Sensitive Groups & Baja & Media & Unhealthy for Sensitive Groups \\
Unhealthy for Sensitive Groups & Baja & Alta & Unhealthy \\
Unhealthy for Sensitive Groups & Media & Baja & Unhealthy for Sensitive Groups \\
Unhealthy for Sensitive Groups & Media & Media & Unhealthy \\
Unhealthy for Sensitive Groups & Media & Alta & Unhealthy \\
Unhealthy for Sensitive Groups & Alta & Baja & Unhealthy \\
Unhealthy for Sensitive Groups & Alta & Media & Unhealthy \\
Unhealthy for Sensitive Groups & Alta & Alta & Very Unhealthy \\
Unhealthy & Baja & Baja & Unhealthy \\
Unhealthy & Baja & Media & Unhealthy \\
Unhealthy & Baja & Alta & Very Unhealthy \\
Unhealthy & Media & Baja & Unhealthy \\
Unhealthy & Media & Media & Very Unhealthy \\
Unhealthy & Media & Alta & Very Unhealthy \\
Unhealthy & Alta & Baja & Very Unhealthy \\
Unhealthy & Alta & Media & Very Unhealthy \\
Unhealthy & Alta & Alta & Hazardous \\
Very Unhealthy & Baja & Baja & Very Unhealthy \\
Very Unhealthy & Baja & Media & Very Unhealthy \\
Very Unhealthy & Baja & Alta & Hazardous \\
Very Unhealthy & Media & Baja & Very Unhealthy \\
Very Unhealthy & Media & Media & Hazardous \\
Very Unhealthy & Media & Alta & Hazardous \\
Very Unhealthy & Alta & Baja & Hazardous \\
Very Unhealthy & Alta & Media & Hazardous \\
Very Unhealthy & Alta & Alta & Hazardous \\
Hazardous & Baja & Baja & Hazardous \\
Hazardous & Baja & Media & Hazardous \\
Hazardous & Baja & Alta & Hazardous \\
Hazardous & Media & Baja & Hazardous \\
Hazardous & Media & Media & Hazardous \\
Hazardous & Media & Alta & Hazardous \\
Hazardous & Alta & Baja & Hazardous \\
Hazardous & Alta & Media & Hazardous \\
Hazardous & Alta & Alta & Hazardous \\
\bottomrule
\end{longtblr}


\clearpage
\VIUAnexoSection{Anexo V}{ane:anexo_breakpoints_epa}
\subsection*{AQI Breakpoints EPA/AQS utilizados por el prototipo}
\begin{longtblr}[
  label={tab:anexo_breakpoints_epa_aqs},
  caption={AQI Breakpoints EPA/AQS utilizados por el prototipo para construir subíndices por contaminante. La tabla resume únicamente los tramos implementados en \textit{AQRisk}.}
]{
  colspec={X[1.5,l,m] X[0.9,l,m] X[1.6,l,m] X[0.7,c,m] X[0.7,c,m] X[0.9,c,m] X[0.9,c,m]},
  width=\linewidth
}
\toprule
Parámetro & Duración & Categoría AQI & AQI bajo & AQI alto & Breakpoint bajo & Breakpoint alto \\
\midrule
PM$_{2.5}$ & 24 h & Good & 0 & 50 & 0.0 & 9.0 \\
PM$_{2.5}$ & 24 h & Moderate & 51 & 100 & 9.1 & 35.4 \\
PM$_{2.5}$ & 24 h & Unhealthy for Sensitive Groups & 101 & 150 & 35.5 & 55.4 \\
PM$_{2.5}$ & 24 h & Unhealthy & 151 & 200 & 55.5 & 125.4 \\
PM$_{2.5}$ & 24 h & Very Unhealthy & 201 & 300 & 125.5 & 225.4 \\
PM$_{2.5}$ & 24 h & Hazardous & 301 & 500 & 225.5 & 325.4 \\
PM$_{10}$ & 24 h & Good & 0 & 50 & 0.0 & 54.0 \\
PM$_{10}$ & 24 h & Moderate & 51 & 100 & 55.0 & 154.0 \\
PM$_{10}$ & 24 h & Unhealthy for Sensitive Groups & 101 & 150 & 155.0 & 254.0 \\
PM$_{10}$ & 24 h & Unhealthy & 151 & 200 & 255.0 & 354.0 \\
PM$_{10}$ & 24 h & Very Unhealthy & 201 & 300 & 355.0 & 424.0 \\
PM$_{10}$ & 24 h & Hazardous & 301 & 500 & 425.0 & 604.0 \\
CO & 8 h & Good & 0 & 50 & 0.0 & 4.4 \\
CO & 8 h & Moderate & 51 & 100 & 4.5 & 9.4 \\
CO & 8 h & Unhealthy for Sensitive Groups & 101 & 150 & 9.5 & 12.4 \\
CO & 8 h & Unhealthy & 151 & 200 & 12.5 & 15.4 \\
CO & 8 h & Very Unhealthy & 201 & 300 & 15.5 & 30.4 \\
CO & 8 h & Hazardous & 301 & 500 & 30.5 & 50.4 \\
NO$_2$ & 1 h & Good & 0 & 50 & 0.0 & 53.0 \\
NO$_2$ & 1 h & Moderate & 51 & 100 & 54.0 & 100.0 \\
NO$_2$ & 1 h & Unhealthy for Sensitive Groups & 101 & 150 & 101.0 & 360.0 \\
NO$_2$ & 1 h & Unhealthy & 151 & 200 & 361.0 & 649.0 \\
NO$_2$ & 1 h & Very Unhealthy & 201 & 300 & 650.0 & 1249.0 \\
NO$_2$ & 1 h & Hazardous & 301 & 500 & 1250.0 & 2049.0 \\
O$_3$ & 8 h & Good & 0 & 50 & 0.000 & 0.054 \\
O$_3$ & 8 h & Moderate & 51 & 100 & 0.055 & 0.070 \\
O$_3$ & 8 h & Unhealthy for Sensitive Groups & 101 & 150 & 0.071 & 0.085 \\
O$_3$ & 8 h & Unhealthy & 151 & 200 & 0.086 & 0.105 \\
O$_3$ & 8 h & Very Unhealthy & 201 & 300 & 0.106 & 0.200 \\
O$_3$ & 1 h & Unhealthy for Sensitive Groups & 101 & 150 & 0.125 & 0.164 \\
O$_3$ & 1 h & Unhealthy & 151 & 200 & 0.165 & 0.204 \\
O$_3$ & 1 h & Very Unhealthy & 201 & 300 & 0.205 & 0.404 \\
O$_3$ & 1 h & Hazardous & 301 & 500 & 0.405 & 0.604 \\
SO$_2$ & 1 h & Good & 0 & 50 & 0.0 & 35.0 \\
SO$_2$ & 1 h & Moderate & 51 & 100 & 36.0 & 75.0 \\
SO$_2$ & 1 h & Unhealthy for Sensitive Groups & 101 & 150 & 76.0 & 185.0 \\
SO$_2$ & 1 h & Unhealthy & 151 & 200 & 186.0 & 304.0 \\
SO$_2$ & 24 h & Very Unhealthy & 201 & 300 & 305.0 & 604.0 \\
SO$_2$ & 24 h & Hazardous & 301 & 500 & 605.0 & 1004.0 \\
\bottomrule
\end{longtblr}

\noindent\footnotesize\textit{Fuente:} tabla oficial \textit{AQI Breakpoints} de EPA/AQS \cite{DATA_EPA_AQS_BREAKPOINTS}. \normalsize

